\documentclass[a4paper,11pt]{article}

\usepackage[margin=2.5cm]{geometry}
\usepackage[T1]{fontenc}
\usepackage[utf8]{inputenc}
\usepackage{lmodern}
\usepackage{booktabs}
\usepackage{enumitem}
\usepackage{graphicx}
\usepackage{float}
\usepackage{xcolor}
\usepackage[hidelinks]{hyperref}
\usepackage[style=ieee]{biblatex}
\addbibresource{library.bib}

\definecolor{darkblue}{RGB}{0, 0, 102}
\hypersetup{
    colorlinks=true,
    urlcolor=darkblue,
    linkcolor=darkblue,
    citecolor=darkblue,
}

\title{
    \textbf{DNS-Based Traffic Steering for IPv6-Enabled Edge Microservices}\\[0.5em]
    \large Extended Project Description
}
\author{Joar Heimonen \\ \texttt{contact@joar.me}}
\date{\today}

\begin{document}

\maketitle
\tableofcontents
\newpage

\section{Introduction}
% - What is the problem? Proxies are everywhere, they introduce overhead and single points of failure
% - Motivate from PEAK: IPv6 abundance eliminates NAT, DNS can steer traffic directly
% - State the thesis question clearly
% - Briefly describe the structure of this document
The modern cloud computing stack relies heavily on Layer 7 proxies and service meshes to manage traffic between services however
this comes at a cost. These layer 7 proxies introduce single points of failure in otherwise distributed systems, whilst adding unneccesary
complexity and reducing the effectiveness of layer 3 routing. Service meshes add sidecar containers to every pod.
This overhead scales with the amount of microservices you have.

This architectural pattern did not emerge by design. IPv4 exhaustion forced the development of technologies like
NAT (Network Address Translation) and shared address spaces. With the introduction of NAT end-to-end accessibility was broken
making proxies a structural necessity.

With a rapid increase in IPv6 adoption it is time to leave the technical limitations of the past in the past
and build our services with IPv6s large adress space in mind. Under RIPE-738 \cite{IPv6AddressAllocationa}, a Local internet registry recieves a minimum allocation size of CIDR /32 up to /29
without needing to justify it. A /29 allocation alone contains around 147 billion times more addresses than the whole IPv4 address space.

This thesis investigates whether DNS-based traffic steering, combined with direct IPv6 addressing of microservices can serve
as a viable alternative to traditional layer 7 ingress proxies and service meshes. Rather than routing traffic trough a centralized proxy
each pod is assigned a globally routable address and made directly reachable by clients.
Traffic steering decissions are delegated to a purpose built DNS server that dynamically manages record
based on pod health and metrics.

The remainder of this document surveys the state of the art in traffic steering and proxy scalability, identifies the gap that motivates this research
and outlines the proposed approach, experimental methodology and work plan.

\section{State of the Art in Traffic Steering and Proxy Scalability}

\subsection{The End-to-End Principle and the Role of NAT}
% Saltzer Reed Clark 1984 — functions belong at endpoints
% IPv4 NAT broke end-to-end addressability
% IPv6 restores it — this is the architectural enabler of PEAK

Saltzer, Reed and Clark established in 1984 that application specific functions should be
implemented at the endpoints of a communication system rather than within the network itself \cite{saltzerEndtoendArgumentsSystem1984}.
A network that simply moves packets without enforcing control points is inherently more ressilient. There are no forced 
bottlenecks, no single points of failure introduces by the infrastructure itself. Packets are free to be
routed as the network sees fit.

The widespread adoption of NAT under IPv4 fundementally broke this principle. By placing address translation logic
inside the network, NAT made end-to-end addressability impossible. Applications could no longer communicate directly over the internet.
They required proxies, STUN servers and other middleboxes to facilitate communication compansating for the loss of
globally unique addressing.

IPv6 restores the conditions under which end-to-end addressability principle can be applied. With a globally unique address assigned to
every endpoint there is no longer a structural requirement for address translation or centralized proxies. The network can return to its intended role,
moving packets between endpoints that are directly reachable by design.



\subsection{Layer 7 Proxies and Ingress Controllers}
% What is a proxy, what does it do
% HAProxy, NGINX, Envoy, Traefik — design patterns and tradeoffs
% Overhead: TLS termination, connection management, CPU cost
% Centralized bottleneck problem
% Cite: bachelor comparison, secure reverse proxy

\subsection{Service Mesh Architectures}
% Why service meshes emerged — observability, resilience, policy
% Sidecar model: Istio, Cilium, AWS App Mesh
% Overhead analysis — duplicate protocol processing, cross-process cost
% Movement toward sidecar-less approaches
% Cite: selecting service mesh, XLB motivation section

\subsection{Kernel and Hardware Acceleration}
% The insight: placement of logic in the stack drives overhead
% eBPF: moving L7 logic into kernel (XLB, Cilium/Miziński)
% DPU: moving service mesh to hardware (FlatProxy)
% Tradeoffs: complexity, compatibility, security model
% Cite: XLB, FlatProxy, eBPF Kubernetes paper

\subsection{DNS-Based Traffic Steering}
% DNS as a distributed system — what it can and cannot do
% GeoDNS: geographic steering at Internet scale
% TTL as the core tradeoff: consistency vs load
% Cite: Hawley, Moura, Shaikh, Chatterjee
% The gap: nobody has used DNS as a primary per-pod steering mechanism
% IPv6 is what makes this viable — globally routable pod addresses

\subsection{Edge and Orchestration Context}
% Kubernetes at the edge — etcd bottleneck, consistency tradeoffs
% Pressure on centralized control assumptions
% Cite: Rearchitecting Kubernetes for the Edge

\section{Literature Search Methodology}
% How the search was conducted
% Keywords used
% Inclusion/exclusion criteria
% Snowballing
% Table of papers

The literature search was conducted using Google Scholar and Google Scholar Labs.
Search terms were initially generated with the assistance of a large language model based on the research topic, and subsequently refined through iterative searching.
The following keyword queries were used:
\begin{itemize}
    \item \textit{DNS TTL load balancing consistency tradeoff}
    \item \textit{IPv6 end-to-end addressing NAT elimination}
    \item \textit{eBPF Kubernetes networking performance}
    \item \textit{service mesh overhead scalability microservices}
    \item \textit{DNS-based server selection load balancing}
\end{itemize}
Papers were selected based on relevance to the research problem, publication venue quality, and citation count.
The search was supplemented by backward snowballing from the reference lists of key papers, through which several foundational works were identified.
Papers published in predatory or low-quality journals were excluded.

The resulting set of nine papers is summarized in Table \ref{tab:papers}.

\begin{table}[H]
\centering
\small
\begin{tabular}{lllp{6cm}}
\toprule
\textbf{\#} & \textbf{Authors} & \textbf{Year} & \textbf{Key Claim} \\
\midrule
1 & Wang, Shou, Qian, Liu          & 2026 & In-kernel eBPF L7 load balancer, 1.5x throughput over Istio \cite{wangXLBHighPerformance2026} \\
2 & Li, Lu, Lin, Wu, Zhang, Yan    & 2023 & DPU-centric service mesh, 90\% latency reduction vs Envoy \cite{liFlatProxyDPUcentricService2023} \\
3 & Hawley                         & 2009 & GeoDNS geographic traffic steering at Internet scale \cite{hawleyGeoDNSGeographicallyAware} \\
4 & Jeffery, Howard, Mortier       & 2021 & etcd strong consistency bottleneck in edge Kubernetes \cite{jefferyRearchitectingKubernetesEdge2021} \\
5 & Miziński, Przyłucki            & 2025 & Cilium eBPF vs iptables: lower CPU/memory under load \cite{mizinskiImpactUsingEBPF2025} \\
6 & Moura, Heidemann, Schmidt      & 2019 & DNS TTL tradeoffs between caching efficiency and routing flexibility \cite{mouraCacheMeIf2019} \\
7 & Shaikh, Tewari, Agrawal        & 2001 & Reducing DNS TTL increases nameserver load by orders of magnitude \cite{shaikhEffectivenessDNSbasedServer2001} \\
8 & Chatterjee, Tari, Zomaya       & 2005 & Adaptive TTL approach reduces client-perceived latency by 16\% \cite{chatterjeeTaskbasedAdaptiveTTL2005} \\
9 & Saltzer, Reed, Clark           & 1984 & End-to-end principle: functions belong at endpoints, not in the network \cite{saltzerEndtoendArgumentsSystem1984} \\
\bottomrule
\end{tabular}
\caption{Papers included in the literature review}
\label{tab:papers}
\end{table}

\section{Problem Description}
% Detailed statement of the research problem
% What gap does this thesis address
% Concrete examples of the limitations of current approaches
% How PEAK (prior work) motivates the research question
% What we do not yet know — the experimental gap

\section{Discussion}
% Comparison across approaches
% Common patterns: all current approaches assume centralized L7 control
% The overhead is not incidental — it is structural
% IPv6 + DNS as a structural alternative
% What tradeoffs does the proposed approach introduce
% Consistency guarantees, TTL propagation, failure detection

\section{Proposed Approach and Methodology}
% Describe the PEAK architecture and PeakDNS
% What the thesis will build on top of the prior work
% Experimental design:
%   - DNS server as Kubernetes ingress controller
%   - IPv6 globally routable pod addresses
%   - Prometheus-based metric-driven record management
%   - Dummy microservices and client simulators
%   - Constant pod count, varying client count
%   - Metrics: RTT, throughput, CPU, failover behavior

\section{Work Plan and Milestones}

\subsection*{Spring 2026}
\begin{itemize}
    \item Systematization of knowledge and literature study
    \item Write and submit mandatory essay on the state of proxy scalability
    \item Restructure essay into draft systematization of knowledge chapter
    \item Start development of DNS server
\end{itemize}

\subsection*{Autumn 2026}
\begin{itemize}
    \item Finish development of DNS server and custom ingress controller
    \item Implement client simulator and dummy microservices
    \item Select baseline paradigms
    \item Start developing the experiments
    \item Begin writing thesis background and methodology chapters
\end{itemize}

\subsection*{Spring 2027}
\begin{itemize}
    \item Perform the experiments
    \item Analyze the results
    \item Write article on DNS-based traffic steering
    \item Finish master thesis
\end{itemize}

\section{Conclusion}
% Summarize the state of the art
% What problems remain unsolved
% What the literature is moving toward
% Where your thesis fits and what it contributes

\newpage
\printbibliography
\end{document}
